\documentclass{article}

\usepackage{fancyhdr}
\usepackage{extramarks}

%amsmath
\usepackage{amsmath}
\usepackage{amsthm}
\usepackage{amsfonts}
\usepackage{mathtools}
\usepackage{amssymb}

%Enumerating items and adding columns
\usepackage{enumitem}
\usepackage{multicol}

%TikZ picture
\usepackage{tikz}
\usepackage{scalerel}
\usepackage{pict2e}
\usepackage{tkz-euclide}
\usetikzlibrary{calc}
\usetikzlibrary{patterns,arrows.meta}
\usetikzlibrary{shadows}
\usetikzlibrary{external}
\usetikzlibrary{automata,positioning}

%pgfplots
\usepackage{pgfplots}
\pgfplotsset{compat=newest}
\usepgfplotslibrary{statistics}
\usepgfplotslibrary{fillbetween}

%colours
\usepackage{xcolor}

%
% Basic Document Settings
%

\topmargin=-0.45in
\evensidemargin=0in
\oddsidemargin=0in
\textwidth=6.5in
\textheight=9.0in
\headsep=0.25in

\linespread{1.1}

\pagestyle{fancy}
\lhead{\hmwkAuthorName}
\chead{\hmwkChapter . \hmwkChapterTitle \hmwkClassTime: \hmwkTitle}
\rhead{\firstxmark}
\lfoot{\lastxmark}
\cfoot{\thepage}

\renewcommand\headrulewidth{0.4pt}
\renewcommand\footrulewidth{0.4pt}

\setlength\parindent{0pt}

%
% Create Problem Sections
%

\newcommand{\enterProblemHeader}[1]{
    \nobreak\extramarks{}{Problem \arabic{#1} continued on next page\ldots}\nobreak{}
    \nobreak\extramarks{Problem \arabic{#1} (continued)}{Problem \arabic{#1} continued on next page\ldots}\nobreak{}
}

\newcommand{\exitProblemHeader}[1]{
    \nobreak\extramarks{Problem \arabic{#1} (continued)}{Problem \arabic{#1} continued on next page\ldots}\nobreak{}
    \stepcounter{#1}
    \nobreak\extramarks{Problem \arabic{#1}}{}\nobreak{}
}

\setcounter{secnumdepth}{0}
\newcounter{partCounter}
\newcounter{homeworkProblemCounter}
\setcounter{homeworkProblemCounter}{1}
\nobreak\extramarks{Problem \arabic{homeworkProblemCounter}}{}\nobreak{}

%
% Homework Problem Environment
%
% This environment takes an optional argument. When given, it will adjust the
% problem counter. This is useful for when the problems given for your
% assignment aren't sequential. See the last 3 problems of this template for an
% example.
%
\newenvironment{homeworkProblem}[1][-1]{
    \ifnum#1>0
        \setcounter{homeworkProblemCounter}{#1}
    \fi
    \section{Problem \arabic{homeworkProblemCounter}}
    \setcounter{partCounter}{1}
    \enterProblemHeader{homeworkProblemCounter}
}{
    \exitProblemHeader{homeworkProblemCounter}
}

%
% Homework Details
%   - Title
%   - Due date
%   - Class
%   - Section/Time
%   - Instructor
%   - Author
%

\newcommand{\hmwkTitle}{Exercises}
\newcommand{\hmwkChapter}{CHAPTER 4}
\newcommand{\hmwkDate}{April 4, 2026}
\newcommand{\hmwkClassTime}{}
\newcommand{\hmwkChapterTitle}{Induction}
\newcommand{\hmwkAuthorName}{\textbf{Christian Jelo R. Artoza}}

%
% Title Page
%

\title{
    \vspace{2in}
    \textmd{\textbf{\hmwkChapter:\ \hmwkTitle}}\\
    \normalsize\vspace{0.1in}\small{\hmwkDate}\\
    \vspace{0.1in}\large{\textit{\hmwkChapterTitle\ \hmwkClassTime}}
    \vspace{3in}
}

\author{\hmwkAuthorName}
\date{}

\renewcommand{\part}[1]{\textbf{\large Part \Alph{partCounter}}\stepcounter{partCounter}\\}

%
% Various Helper Commands
%

% qedsymbol that is always flushed to the right
\newcommand{\qedright}{
    \begin{flushright}
        \qedsymbol
    \end{flushright}
}

% For propositions
\newtheorem{prop}{Proposition}

% Useful for algorithms
\newcommand{\alg}[1]{\textsc{\bfseries \footnotesize #1}}

% For derivatives
\newcommand{\deriv}[1]{\frac{\mathrm{d}}{\mathrm{d}x} (#1)}

% For partial derivatives
\newcommand{\pderiv}[2]{\frac{\partial}{\partial #1} (#2)}

% Integral dx
\newcommand{\dx}{\mathrm{d}x}

% Alias for the Solution section header
\newcommand{\solution}{\textbf{\large Solution}}

% Probability commands: Expectation, Variance, Covariance, Bias
\newcommand{\E}{\mathrm{E}}
\newcommand{\Var}{\mathrm{Var}}
\newcommand{\Cov}{\mathrm{Cov}}
\newcommand{\Bias}{\mathrm{Bias}}

\begin{document}

\maketitle

\pagebreak

%Plot Settings
\pgfplotsset{
    standard/.style={
    axis line style = thick,
    axis equal,
    trig format=rad,
    enlargelimits,
    axis x line=middle,
    axis y line=middle,
    axis line style={latex-latex},
    enlarge x limits=0.15,
    enlarge y limits=0.15,
    every axis x label/.style={at={(current axis.right of origin)},anchor=north west},
    every axis y label/.style={at={(current axis.above origin)},anchor=south east}
    }
}

\begin{homeworkProblem}
    \textit{(Exercise 4.2.)} Provide these proofs that if $n\in\mathbb{N}$, then $n^2-n$ is even.
    \begin{enumerate}[label=(\alph*)]
        \item Prove it by cases, by considering the "$n$ is even" and "$n$ is odd" cases.
        \item Prove it by applying \textit{Proposition 4.2} to the sum $1+2+3+\cdots+(n-1)$.
        \item Prove it by induction.
    \end{enumerate}
    \begin{enumerate}[label=(\alph*)]
        \item \begin{proof}
            Let $n\in\mathbb{N}$. Then either $n$ is even, or $n$ is odd. We'll consider these
            cases, separately.
            \vspace{1.5mm}

            \underline{Case 1.} $n$ is even. \\
            Suppose $n$ is even. Then, since the square of an even integer is even, it follows that
            $n^2$ is even. Combining with the fact that the difference of two even integers is even,
            we see that $n^2 - n$ is even, as desired.
            \vspace{1.5mm}

            \underline{Case 2.} $n$ is odd. \\
            Suppose $n$ is odd. Then $n-1$ is even. So, since the product of an even integer and
            an odd integer is even, it follows that the product $n(n-1)$ is even. But, observe that
            $n^2-n = n(n-1)$. Thus, $n^2-n$ is even.
            \vspace{1.5mm}

            We have shown that $n^2-n$ is even whether $n$ is odd, or even. This completes the proof.
            \end{proof}
        \item \begin{proof}
            Let $n\in\mathbb{N}$. Then, we can apply \textit{Proposition 4.2} on $n-1$, and get
            \begin{align*}
                1+2+3+\cdots+(n-1)=\frac{n(n-1)}{2}
            \end{align*}
            Multiplying $2$ to both sides, we get
            \begin{align*}
                2(1+2+3+\cdots+(n-1))=n(n-1).
            \end{align*}
            Thus,
            \begin{align*}
                n^2-n=2(1+2+3+\cdots+(n-1)).
            \end{align*}
            Since the sum at the right hand side of the equation is an integer, it follows that
            $n^2-n$ is an even number, as desired.
            \end{proof} 
        \item \begin{proof}
            We proceed by induction.
            \vspace{1.5mm}

            \underline{Base Case}. The base case is when $n=1$, and
            \begin{align*}
                1^2 - 1 = 0,
            \end{align*}
            which is an even number, as desired.
            \vspace{1.5mm}

            \underline{Inductive Hypothesis}. Suppose that the result holds for all
            $k\in \mathbb{N}$. That is, for any $k\in\mathbb{N}$, $k^2-k$ is even.
            \vspace{1.5mm}

            \underline{Induction Step}. We want to show that the result holds for $k+1$. This means
            proving that $(k+1)^2-(k+1)$ is even. First, observe that
            \begin{align*}
                (k+1)^2-(k+1) &= (k^2+2k+1)-(k+1) \\
                &= (k^2-k) + 2k.
            \end{align*}
            By the inductive hypothesis, the first term is even; and the second term is clearly
            an even number as well. So, since the sum of two even numbers is even, it follows that
            $(k+1)^2-(k+1)$ is even.
            \vspace{1.5mm}

            \underline{Conclusion}. Therefore, the result holds for all $n\in\mathbb{N}$.
            \end{proof}
    \end{enumerate}
\end{homeworkProblem}

\begin{homeworkProblem}
    \textit{(Exercise 4.3.)} Use induction or strong induction to prove that the following hold
    for every $n\in \mathbb{N}$.
    \begin{multicols}{3}
        \begin{enumerate}[label=(\alph*)]
            \item $3\mid (4^n-1)$
            \item $6\mid (n^3-n)$
            \columnbreak
            \item $9\mid (3^{4n}+9)$
            \item $5\mid (n^5-n)$
            \columnbreak
            \item $6\mid (5^{2n}-1)$
            \item $5\mid (6^n - 1)$
        \end{enumerate}
    \end{multicols}
    \textbf{Solution.}
    \begin{enumerate}[label=(\alph*)]
        \item \begin{proof}
            The base case is when $n=1$, and $4^1-1=3$, so $3\mid (4^1-1)$. Now, suppose for any
            $k\in\mathbb{N}$, $3\mid (4^k-1)$. We wish to show that the result holds for $k+1$.
            Observe that
            \begin{align*}
                4^{(k+1)}-1 &= 4\cdot 4^{k}-1 \\
                &= 3\cdot 4^{k} + 4^k-1.
            \end{align*}
            The first term is a multiple of $3$, and hence is divisible by $3$. And, by the induction
            hypothesis, the second term is also divisible by $3$. Hence, the sum must also be divisible
            by $3$, as desired. Therefore, by induction, $3\mid (4^n-1)$ for all $n\in\mathbb{N}$.
        \end{proof}
        \item \begin{proof}
            For the base case, we have $1^3-1=0$. And, indeed, $6\mid 0$. Now, suppose for any
            $k\in \mathbb{N}$, $6\mid(k^3-k)$. To show that the result holds for $k+1$, note that
            we can write
            \begin{align*}
                (k+1)^3-(k+1) &= (k^3+3k^2+3k+1)-(k+1) \\
                &= (k^3-k)+3k(k+1)
            \end{align*}
            By the inductive hypothesis, the first term is divisible by $6$. And, for the second term,
            note that either $k$ is even or $k+1$ is even. Without loss of generality, suppose $k$
            is even. Then, $k=2\ell$ for some $\ell \in \mathbb{N}$. Thus,
            \begin{align*}
                3k(k+1) &= 3\cdot 2\ell\cdot (2\ell+1)\\
                &=6\ell(2\ell+1)
            \end{align*}
            Hence, the second term is also divisible by $6$, and thus $6\mid ((k+1)^3-(k+1))$. We have
            shown that, for any $k\in \mathbb{N}$, $6\mid (k^3-k)$ implies $6\mid ((k+1)^3-(k+1))$. Therefore,
            by induction, $6\mid (n^3-n)$ for all $n\in\mathbb{N}$.
        \end{proof}
        \item \begin{proof}
            The base case is when $n=1$, and $3^{4\cdot 1}+9=90$. Hence, $9\mid 3^{4\cdot 1}+9$,
            as desired. Now, suppose the result holds for any $k\in\mathbb{N}$. To prove the $k+1$
            case, observe that we can write
            \begin{align*}
                3^{4(k+1)}+9 &= 3^4\cdot3^{4k}+9 \\
                &= (3^4-1)\cdot 3^{4k} + (3^{4k}+9) \\
                &= 9^{2k}\cdot(3^4-1) + (3^{4k}+9).
            \end{align*}
            Since $2k\geq 2$, it follows that the first term is divisible by $9$. Furthermore,
            it follows by the inductive hypotheses that the second term is divisible by $9$.
            Thus, $9\mid (3^{4(k+1)}+9)$. Therefore, the result must hold for all $n\in\mathbb{N}$.
        \end{proof}
        \item \begin{proof}
            The base case is when $n=1$, and $1^5-1=0$. Hence, $5\mid 1^5 -1$, as desired. Now, 
            assume that $5\mid (k^5-k)$ for any $k\in \mathbb{N}$. We want to prove that 
            $5\mid ((k+1)^5-(k+1))$. To accomplish this, observe that
            \begin{align*}
                (k+1)^5-(k+1) &= (k^5+5k^4+10k^3+10k^2+5k+1)-(k+1) \\
                &= (k^5-k)+(5k^4+10k^3+10k^2+5k) \\
                &= (k^5-k)+5\cdot (k^4+2k^3+2k^2+k)
            \end{align*}
            By the inductive hypothesis, the first term is divisible by $5$; and the second term
            is also divisible by $5$. Hence, the result holds for $k+1$. Therefore, it must hold
            for all $n\in\mathbb{N}$.
        \end{proof}
        \item \begin{proof}
            Since $5^(2\cdot 1)-1=24$ and $6\mid 24$, the result holds for the base case, $n=1$.
            Now, suppose the result holds for any $k\in\mathbb{N}$. We wish to show that
            $6\mid 5^{2(k+1)}-1$. We have
            \begin{align*}
                5^{2(k+1)}-1 &= 5^2 \cdot 5^{2k}-1 \\
                &= 24\cdot 5^{2k} + (5^{2k}-1) \\
                &= 6\cdot 4\cdot 5^{2k} + (5^{2k}-1)
            \end{align*}
            The first term is a multiple of $6$, so it is divisible by $6$. And, by the induction
            hypothesis, the second term is divisible by $6$. Thus, $6\mid 5^{2(k+1)}-1$ as desired.
            Therefore, the result must hold for all $n\in\mathbb{N}$.
        \end{proof}
        \item \begin{proof}
            For the base case, we know that $6^0 -1=0$ and $5\mid 0$. So, the result holds for the
            base case. Now suppose $5\mid (6^k-1)$ for any $k\in\mathbb{N}$. We want to prove that
            the result also holds for $k+1$. To this end, observe that
            \begin{align*}
                6^{k+1}-1 &= 6\cdot 6^k -1\\
                &= 5\cdot 6^k + (6^k-1)
            \end{align*}
            The first term is divisible by $5$, and, by the inductive hypothesis, the second term
            is divisible by $5$. Hence, $5\mid (6^{k+1}-1)$, and therefore $5\mid (6^n-1)$ for 
            all $n\in \mathbb{N}$.
        \end{proof}
    \end{enumerate}
\end{homeworkProblem}

\begin{homeworkProblem}
    \textit{(Exercise 4.4.)} Prove that each of the following hold for every $n\in\mathbb{N}$.
    \begin{enumerate}[label=(\alph*)]
        \item $1^2+2^2+3^2+\cdots+n^2=\frac{n(n+1)(2n+1)}{6}$
        \item $1^3+2^3+3^3+\cdots+n^3=\frac{n^2(n+1)^2}{4}$
        \item $1\cdot2+2\cdot3+3\cdot4+\cdots+n\cdot(n+1)=\frac{n(n+1)(n+2)}{3}$
        \item $1\cdot3+2\cdot4+3\cdot5+\cdots+n\cdot(n+2)=\frac{n(n+1)(2n+7)}{6}$
        \item $1^3+2^3+3^3+\cdots+n^3=(1+2+3+\cdots+n)^2$
        \item $1\cdot1!+2\cdot2!+3\cdot3!+\cdots+n\cdot n!=(n+1)!-1$
        \item $2^0+2^1+2^2+2^3+\cdots+2^n=2^{n+1}-1$
        \item $3^1+3^2+3^3+\cdots+3^n=\frac{3^{n+1}-3}{2}$
        \item $4^0+4^1+4^2+4^3+\cdots+4^n\frac{4^{n+1}-1}{3}$
        \item $\frac{1}{2!}+\frac{2}{3!}+\frac{3}{4!}+\cdots+\frac{n}{(n+1)!}=1-\frac{1}{(n+1)!}$
        \item $\frac{1}{1\cdot2}+\frac{1}{2\cdot3}+\frac{1}{3\cdot4}+\cdots+\frac{1}{n\cdot(n+1)}
        =\frac{n}{n+1}$
    \end{enumerate}
    \textbf{Solution.}
    \begin{enumerate}[label=(\alph*)]
        \item \begin{proof}
            The base case is when $n=1$, and
            \begin{align*}
                \frac{1\cdot(1+1)(2\cdot1+1)}{6}=\frac{6}{6}=1=1^2,
            \end{align*}
            as desired. Now, suppose the result holds for any $k\in\mathbb{N}$. We want to show
            that the result also holds for $k+1$. That is, we want to prove that
            \begin{align*}
                1^2+2^2+3^2+\cdots+k^2+(k+1)^2=\frac{(k+1)((k+1)+1)(2(k+1)+1)}{6}.
            \end{align*}
            Equivalently, we wish to show that
            \begin{align*}
                1^2+2^2+3^2+\cdots+k^2+(k+1)^2=\frac{(k+1)(k+2)(2k+3)}{6}.
            \end{align*}
            We begin by manipulating the left hand side of the equation. By the inductive
            hypothesis, we have
            \begin{align*}
                1^2+2^2+3^2+\cdots+k^2+(k+1)^2 &= \frac{k(k+1)(2k+1)}{6}+(k+1)^2 \\
                &= \frac{k(k+1)(2k+1)+6(k+1)^2}{6} \\
                &= \frac{(k+1)(2k^2+k+6k+6)}{6} \\
                &= \frac{(k+1)(2k^2+4k+3k+6)}{6} \\
                &= \frac{(k+1)(k+2)(2k+3)}{6},
            \end{align*}
            which is what we wanted. So, the result also holds for $k+1$, and therefore it must
            hold for all $n\in \mathbb{N}$.
        \end{proof}
        \item \begin{proof}
            For the base case, we have,
            \begin{align*}
                \frac{1^2\cdot(1+1)^2}{4} = 1 = 1^3,
            \end{align*}
            as desired. Now, assume that for any $k\in\mathbb{N}$,
            \begin{align*}
                1^3+2^3+3^3+\cdots+k^3=\frac{k^2(k+1)^2}{4}.
            \end{align*}
            We wish to show that the result also holds for $k+1$. In other words, we want to
            prove that
            \begin{align*}
                1^3+2^3+3^3+\cdots+k^3+(k+1)^3=\frac{(k+1)^2((k+1)+1)^2}{4},
            \end{align*}
            or, equivalently, that
            \begin{align*}
                1^3+2^3+3^3+\cdots+k^3+(k+1)^3=\frac{(k+1)^2((k+2)^2}{4}.
            \end{align*}
            We begin with the left hand side of the equation. Using the inductive hypothesis,
            coupled with some algebraic manipulations, we see that
            \begin{align*}
                1^3+2^3+3^3+\cdots+k^3+(k+1)^3 &=\frac{k^2(k+1)^2}{4}+(k+1)^3, \\
                &= \frac{(k+1)^2(k^2+4k+4)}{4},\\
                &= \frac{(k+1)^2(k+2)^2}{4},
            \end{align*}
            as desired. Therefore, the result must hold for all $n\in\mathbb{N}$
        \end{proof}
        \item \begin{proof}
            For $n=1$, the result states that
            \begin{align*}
                1\cdot 2 = \frac{1\cdot(1+1)\cdot(1+2)}{3},
            \end{align*}
            which is clearly true. Now, suppose the result holds for any $k\in\mathbb{N}$. Then,
            \begin{align*}
                1\cdot2+2\cdot3+3\cdot4+\cdots+n\cdot(n+1)+(n+1)\cdot(n+2)
                &=\frac{n(n+1)(n+2)}{3}+(n+1)\cdot(n+2),\\
                &=\frac{(n+1)(n+2)(n+3)}{3},\\
                &=\frac{(n+1)((n+1)+1)((n+1)+2)}{3}.
            \end{align*}
            Hence, the result also holds for $k+1$. Therefore, it must hold for all $n\in\mathbb{N}$.
        \end{proof}
        \item \begin{proof}
            When $n=1$, we have
        \end{proof}
    \end{enumerate}
\end{homeworkProblem}
\end{document}